\section{Ordered products, cyclic indices and path sums}

Throughout this section, ordered products are noncommutative and their
factors are indexed by $\Fin{L}$ in increasing order.

\begin{lemma}[Ordered products of sums and of scalar multiples]
    \label{lem:elementary_ordered_product_sum}
    \lean{List.prod_ofFn_sum, List.prod_ofFn_smul, Fin.prod_succ',
        Fin.prod_mul_mul_eq, Fintype.sum_fin_succ_pi}
    \leanok
    Let $R$ be a semiring, $J$ a finite set and $f_n(j)\in R$ for
    $0\le n<N$ and $j\in J$. Then
    \begin{align}
        \prod_{n=0}^{N-1}\sum_{j\in J}f_n(j)
         & =\sum_{c\colon\Fin{N}\to J}\prod_{n=0}^{N-1}f_n(c(n)),
        \notag
    \end{align}
    the products being ordered. If $R$ is an algebra over a commutative
    semiring $K$, then
    $\prod_n(c_nA_n)=\bigl(\prod_nc_n\bigr)\prod_nA_n$ for scalars
    $c_n\in K$. Moreover $x_0\cdots x_m=(x_0\cdots x_{m-1})x_m$ in every
    monoid, and for matrices $E$ and $F_0,\ldots,F_{m-1}$ of compatible sizes,
    $(EF_0)\cdots(EF_{m-1})E=E(F_0E)\cdots(F_{m-1}E)$. Finally, a sum over
    tuples $\tau\colon\Fin{m+1}\to\alpha$ of a function of $\tau(0)$ and of the
    remaining tuple splits into a sum over the first value and a sum over the
    remaining tuple.
\end{lemma}

\begin{proof}
    \leanok
    The first two identities follow by induction on $N$, splitting off the
    first factor and identifying a choice function on $\Fin{N+1}$ with its
    first value and the remaining choice function. The third is the last-factor
    form of the ordered product. The fourth is an induction on $m$ using
    associativity. The last is the reindexing of tuples of length $m+1$ by
    their first entry and their tail.
\end{proof}

\begin{lemma}[Nonempty products in a generated subsemigroup]
    \label{lem:subsemigroup_closure_nonempty_product}
    \lean{Subsemigroup.exists_nonempty_list_prod_of_mem_closure}
    \leanok
    Let $M$ be a monoid and let $S\subseteq M$. Every element $x$ of the
    subsemigroup generated by $S$ admits a factorization $x=a_1\cdots a_n$
    with $n\geq1$ and $a_j\in S$ for $1\leq j\leq n$.
\end{lemma}

\begin{proof}
    \leanok
    A generator has the required one-factor expression. The product of two
    such expressions is represented by their concatenation, which is again
    nonempty and has every factor in $S$.
\end{proof}

\begin{lemma}[Appending to a finite tuple]
    \label{lem:elementary_list_of_fin_snoc}
    \lean{List.ofFn_snoc, List.exists_eq_ofFn}
    \leanok
    The list of values of a tuple $(f_0,\ldots,f_{n-1},x)$ is the list of
    values of $(f_0,\ldots,f_{n-1})$ followed by $x$. Every finite list is
    the list of values of a tuple indexed by $\Fin{n}$, where $n$ is its
    length.
\end{lemma}

\begin{proof}
    \leanok
    Split the tuple of length $n+1$ at its last entry. A list is the list
    of values of its own entry function.
\end{proof}

\begin{lemma}[Induction on a finite cyclic index]
    \label{lem:elementary_fin_cyclic_induction}
    \lean{Fin.cyclic_induction, Fin.forall_sub_one_iff,
        Fin.finRotate_succ_castSucc, coe_finRotate_mod}
    \leanok
    Let $m\geq1$ and let $P$ be a property of residues modulo $m$. If $P(0)$
    holds and $P(i)$ implies $P(i+1)$ for every $i$, then $P$ holds for every
    residue. A statement about every pair $(j,j-1)$ is the same as the
    statement about every pair $(j+1,j)$. One-step rotation of $\Fin{N+1}$
    sends $i<N$ to $i+1$, and on $\Fin{N}$ it sends $i$ to $(i+1)\bmod N$.
\end{lemma}

\begin{proof}
    \leanok
    Induct on the value of the residue: a nonzero residue is the successor
    of the residue one smaller. The second statement substitutes $j+1$ for
    $j$, and conversely. The rotation identities are computations of values.
\end{proof}

\begin{lemma}[Constancy on a finite interval]
    \label{lem:finite_interval_constancy}
    \lean{eq_of_forall_succ_eq_of_mem_Icc}
    \leanok
    Let $a,n\in\mathbb N$, and let $f$ be defined on the integer interval
    $[0,n]$. If $f(m)=f(m+1)$ for every $a\le m<n$, then $f(i)=f(j)$ for all
    $i,j\in[a,n]$.
\end{lemma}

\begin{proof}
    \leanok
    Composing the successive equalities from the smaller index to the larger
    gives $f(i)=f(j)$.
\end{proof}

\begin{lemma}[Cyclic path expansion of a trace]
    \label{lem:matrix_cyclic_path_sum}
    \lean{Matrix.sum_fin_cons, Matrix.ofFn_prod_mulVec_apply,
        Matrix.ofFn_prod_apply, Matrix.trace_ofFn_prod_eq_sum_cyclic}
    \leanok
    Let $M_0,\ldots,M_{L-1}$ be square matrices of the same finite size over
    a commutative semiring. Applying the ordered product $M_0\cdots M_{L-1}$
    to a vector $v$ and reading the entry at $a$ sums, over all paths
    $(t_0,\ldots,t_L)$ with $t_0=a$, the product of the entries
    $(M_n)_{t_n t_{n+1}}$ along the path times $v_{t_L}$. In particular the
    entry of the product at $(a,b)$ is the sum over the paths from $a$ to $b$.
    If $L\ge1$, then
    \begin{align}
        \operatorname{tr}(M_0\cdots M_{L-1})
         & =\sum_{t_0,\ldots,t_{L-1}}\prod_{n=0}^{L-1}(M_n)_{t_n t_{n+1}},
        \notag
    \end{align}
    with the cyclic convention $t_L=t_0$.
\end{lemma}

\begin{proof}
    \leanok
    \uses{lem:elementary_fin_cyclic_induction}
    The open expansion is an induction on $L$: splitting a path into its
    first step and the remaining path gives
    $(M_0(M_1\cdots M_{L-1})v)_a=\sum_c(M_0)_{ac}((M_1\cdots M_{L-1})v)_c$.
    The entry formula applies the open expansion to a standard basis vector.
    For the trace, separating the last factor $M_{L-1}$ and applying the
    open expansion to the vector $c\mapsto(M_{L-1})_{ca}$, then splitting
    each cyclic assignment into its first entry $a$ and the remaining
    entries, identifies the two sides term by term as the last factor closes
    the cycle back to $a$.
\end{proof}

\begin{lemma}[Cyclic trace of matrix units]
    \label{lem:asymex_ising_unit_loop}
    \lean{Matrix.trace_ofFn_prod_smul_single}
    \leanok
    \uses{lem:matrix_cyclic_path_sum}
    For $N\geq1$, scalars $c_v$, and matrix-unit indices $a_v,b_v$,
    \begin{align}
        \operatorname{tr}\!\left(\prod_{v=0}^{N-1}c_vE_{a_vb_v}\right)
         & =\left(\prod_{v=0}^{N-1}c_v\right)
        \prod_{v=0}^{N-1}\delta_{b_v,a_{v+1}},
        \label{eq:asymex_ising_unit_loop}
    \end{align}
    where $a_N=a_0$.
\end{lemma}

\begin{proof}
    \leanok
    In the cyclic path expansion only the path $t_v=a_v$ can contribute,
    and it contributes $\prod_vc_v$ exactly when $b_v=a_{v+1}$ for every
    $v$.
\end{proof}

\begin{lemma}[Three consecutive positive power sums]
    \label{thm:positive_power_second_difference}
    \lean{Real.sum_pow_secondDifference,
        Real.eq_one_of_sum_pow_secondDifference_eq_zero,
        Complex.eq_one_of_sum_pow_secondDifference_eq_zero}
    \leanok
    Let $x_0,\ldots,x_{r-1}$ be real numbers and $p_n=\sum_k x_k^n$.
    For every integer $n\geq0$,
    \begin{align}
        p_{n+2}-2p_{n+1}+p_n
         & =\sum_k x_k^n(x_k-1)^2.
        \label{eq:elementary_positive_second_difference}
    \end{align}
    If every $x_k$ is positive, the second difference vanishes if and only
    if every $x_k$ equals one. The same holds for complex numbers that are
    positive real numbers.
\end{lemma}

\begin{proof}
    \leanok
    Expanding the squares gives~(\ref{eq:elementary_positive_second_difference}).
    Every summand is nonnegative, and every factor $x_k^n$ is strictly
    positive. Thus the sum vanishes exactly when all $x_k-1$ vanish.
\end{proof}
